\section{Justificación}

% ===============================
% 1. Problema: reglas hechas a mano
% ===============================
\begin{frame}{Justificación}
	\begin{center}
		\begin{tikzpicture}
			\node[cajaobjetivo, text width=0.85\textwidth]
			{
				\textbf{Limitación actual} \\
				El diseño de la función de transición en autómatas celulares depende de expertos y debe hacerse manualmente, lo cual restringe sus aplicaciones prácticas.
			};
		\end{tikzpicture}
	\end{center}
	
	\begin{center}
		\begin{tikzpicture}
			\node[cajaobjetivo, text width=0.85\textwidth]
			{
				\textbf{Complejidad de los sistemas reales} \\
				Muchos fenómenos naturales, sociales o biológicos no poseen reglas explícitas o son demasiado complejos para modelarse de forma analítica.
			};
		\end{tikzpicture}
	\end{center}
\end{frame}

% ===============================
% 2. Necesidad del proyecto
% ===============================
\begin{frame}{Justificación}
	\begin{center}
		\begin{tikzpicture}
			\node[cajaobjetivo, text width=0.85\textwidth]
			{
				\textbf{Necesidad} \\
				Se requiere un enfoque capaz de aprender comportamientos sin depender de funciones matemáticas definidas manualmente.
			};
		\end{tikzpicture}
	\end{center}
	
	\begin{center}
		\begin{tikzpicture}
			\node[cajaobjetivo, text width=0.85\textwidth]
			{
				\textbf{Nuevo enfoque} \\
				Los autómatas celulares neuronales permiten que la función de transición sea aprendida directamente a partir de datos.
			};
		\end{tikzpicture}
	\end{center}
\end{frame}

% ===============================
% 3. Solución y beneficios
% ===============================
\begin{frame}{Justificación}
	\begin{center}
		\begin{tikzpicture}
			\node[cajaobjetivo, text width=0.85\textwidth]
			{
				\textbf{Aprovechamiento de IA} \\
				Las redes convolucionales y recurrentes pueden capturar dependencias espaciales y temporales, siendo adecuadas para aproximar funciones de transición.
			};
		\end{tikzpicture}
	\end{center}
	
	\begin{center}
		\begin{tikzpicture}
			\node[cajaobjetivo, text width=0.85\textwidth]
			{
				\textbf{Reducción de complejidad} \\
				Permiten simular sistemas complejos sin ajustes manuales, evitando procesos costosos o inviables.
			};
		\end{tikzpicture}
	\end{center}
\end{frame}

% ===============================
% 4. Valor científico y futuro
% ===============================
\begin{frame}{Justificación}
	\begin{center}
		\begin{tikzpicture}
			\node[cajaobjetivo, text width=0.85\textwidth]
			{
				\textbf{Aporte científico} \\
				El proyecto integra inteligencia artificial, vida artificial y simulación, creando una plataforma flexible para futuras investigaciones.
			};
		\end{tikzpicture}
	\end{center}
	
	\begin{center}
		\begin{tikzpicture}
			\node[cajaobjetivo, text width=0.85\textwidth]
			{
				\textbf{Impacto a largo plazo} \\
				Se genera un entorno de código abierto, reutilizable y extensible, que puede impulsar nuevas líneas de trabajo en sistemas complejos e IA.
			};
		\end{tikzpicture}
	\end{center}
\end{frame}
